\section{Tagged interleavings}
\label{sec:tagged_interleavings}

Let $A$ and $B$ be arbitrary alphabets, and distinguish them by the two
inclusions into $A\sqcup B$.  For a word $w$ over $A\sqcup B$, let
$\pi_A(w)$ and $\pi_B(w)$ be the words obtained by retaining, in their
original order, the letters carrying the respective tag and then removing
that tag.  Write $|w|$ for the length of a word.

\begin{definition}[Enumeration of tagged interleavings]
\label{def:tagged_interleavings}
\lean{TaggedInterleavings.taggedInterleavings}
\leanok
For words $u$ over $A$ and $v$ over $B$, define the list $I(u,v)$ recursively.
If either word is empty, the list contains just the other word with its tag
attached.  If $u=au'$ and $v=bv'$ are nonempty, first list the words obtained
by prepending the left-tagged $a$ to each member of $I(u',v)$, then list the
words obtained by prepending the right-tagged $b$ to each member of $I(u,v')$.
\end{definition}

\begin{lemma}[Characterization by filtered words]
\label{lem:tagged_interleavings_filter}
\lean{TaggedInterleavings.mem_taggedInterleavings_iff}
\leanok
\uses{def:tagged_interleavings}
A word $w$ occurs in $I(u,v)$ if and only if
$\pi_A(w)=u$ and $\pi_B(w)=v$.
\end{lemma}
\begin{proof}
\leanok
Induct on the two prescribed words.  When one is empty, every letter must
carry the other tag, so the remaining filtered word uniquely determines
$w$.  When both are nonempty, the first tag of $w$ determines which recursive
case applies.  Removing that first letter reduces the assertion to the
corresponding shorter pair of words.
\end{proof}

\begin{lemma}[Number of tagged interleavings]
\label{lem:tagged_interleavings_count}
\lean{TaggedInterleavings.nodup_taggedInterleavings,
TaggedInterleavings.length_taggedInterleavings}
\leanok
\uses{def:tagged_interleavings}
The list $I(u,v)$ has no repeated words and has
$\binom{|u|+|v|}{|u|}$ members, even if $u$ or $v$ repeats letters.
\end{lemma}
\begin{proof}
\leanok
Induct on the two prescribed words.  Prepending a fixed letter preserves
distinctness, and the two nonempty recursive cases have different first
tags.  Thus each recursive step preserves the absence of repeated words.
If either word is empty, there is one interleaving.  Otherwise the number
of interleavings is the sum of the two shorter counts, which equals
$\binom{|u|+|v|}{|u|}$ by Pascal's identity.
\end{proof}

\begin{lemma}[Lengths of the filtered words]
\label{lem:tagged_interleavings_length}
\lean{TaggedInterleavings.length_filterMap_add,
TaggedInterleavings.length_of_mem_taggedInterleavings}
\leanok
\uses{def:tagged_interleavings}
Every word over $A\sqcup B$ satisfies
$|\pi_A(w)|+|\pi_B(w)|=|w|$.
Consequently, every member of $I(u,v)$ has length $|u|+|v|$.
\end{lemma}
\begin{proof}
\leanok
\uses{lem:tagged_interleavings_filter}
Each letter contributes one to exactly one of the two filtered lengths.
Induction on $w$ gives the first identity.  Substitute the two prescribed
filtered words to obtain the consequence.
\end{proof}
